%%%%%%%%%%%%%%%%%%%%%%%%%%%%%%%%%%%%%%%%%
% Embedded Evaluation: Stata Tools for Applied Research and Evaluation
% Eric Booth & Elizabeth Teas
% 2026 Edition
%%%%%%%%%%%%%%%%%%%%%%%%%%%%%%%%%%%%%%%%%

\documentclass[
    a4paper,
    fontsize=10pt,
    twoside=true,
    numbers=noenddot,
]{kaobook}

% --- Geometry Adjustments (5% wider narrative, narrower margins) ---
\addtolength{\textwidth}{0.05\textwidth}
\addtolength{\marginparwidth}{-0.05\textwidth}

% --- Language & Typography ---
\ifxetexorluatex
    \usepackage{polyglossia}
    \setmainlanguage{english}
\else
    \usepackage[english]{babel}
\fi
\usepackage[english=american]{csquotes}

% --- Essential Packages ---
\usepackage{kaobiblio}
\IfFileExists{main.bib}{\addbibresource{main.bib}}{} 
\usepackage[framed=true]{kaotheorems}
\usepackage{kaorefs}

% --- Fallback Definitions for Custom kaobook Commands ---
\providecommand{\Environment}[1]{\texttt{#1}}
\providecommand{\Command}[1]{\texttt{\textbackslash#1}}
\providecommand{\Option}[1]{\textsf{#1}}

% --- Custom asides: "Applied Example" pull-out + "Visualization" callout ---
% Built on mdframed (already loaded by kao.sty/kaotheorems) so they compile
% with the rest of the book. Applied Example = a standalone, colored, boxed
% deep-dive that breaks the chapter narrative. Visualization = a smaller
% prompt suggesting a chart/figure to build around.
\definecolor{aeBg}{HTML}{FBF1DF}      % warm sand background
\definecolor{aeRule}{HTML}{B5651D}    % terracotta rule
\definecolor{aeTitle}{HTML}{7A3E0E}   % dark brown title text
\definecolor{vizBg}{HTML}{E6F3EF}     % pale teal background
\definecolor{vizRule}{HTML}{14776A}   % deep teal rule

\newcounter{appliedexample}
\counterwithin{appliedexample}{chapter}
\renewcommand*{\theHappliedexample}{\thechapter.\arabic{appliedexample}}

\newmdenv[
    skipabove=1.2\topskip, skipbelow=0.8\topskip,
    innertopmargin=8pt, innerbottommargin=8pt,
    innerleftmargin=11pt, innerrightmargin=11pt,
    linewidth=1.1pt, roundcorner=5pt,
    linecolor=aeRule, backgroundcolor=aeBg,
    frametitlerule=true,
    frametitlefont=\sffamily\bfseries\color{aeTitle},
    frametitlebackgroundcolor=aeBg,
]{appliedexamplebox}

% Applied Examples span the text and margin columns (kaobook's widepar) and
% never split across pages: the body is set in a box at the wide measure
% first; if it fits in the space left on the page it is set in place,
% otherwise it floats to the top of the next page while the text continues. The body is collected as an argument, so an Applied
% Example cannot contain verbatim code (none does).
\newsavebox{\aeBoxContent}
% \aeNoFloat, written just before an Applied Example, makes that example start
% on the next page instead of floating there. Use it where a float pending at
% a page break pushes a following \section* heading onto a new page.
\newif\ifaenofloat
\newcommand{\aeNoFloat}{\global\aenofloattrue}
\NewDocumentEnvironment{appliedexample}{O{} +b}{%
    \par
    \refstepcounter{appliedexample}%
    \sbox{\aeBoxContent}{\begin{minipage}{\dimexpr\textwidth+\marginparsep+\marginparwidth\relax}%
        \begin{appliedexamplebox}[frametitle={Applied Example~\theappliedexample.~~#1},nobreak=true]%
        #2%
        \end{appliedexamplebox}%
    \end{minipage}}%
    \ifdim\dimexpr\pagegoal-\pagetotal<\dimexpr\ht\aeBoxContent+\dp\aeBoxContent+1.2\topskip\relax
        \ifaenofloat
            \newpage
            \begin{widepar}\noindent\usebox{\aeBoxContent}\end{widepar}%
            \par
        \else
            % Not enough room here: float the box to the top of the next
            % page, using the same wide float the book's figure* environments
            % use, and let the following text fill the rest of this page.
            \begin{figure*}[tp]\noindent\usebox{\aeBoxContent}\end{figure*}%
        \fi
    \else
        \begin{widepar}\noindent\usebox{\aeBoxContent}\end{widepar}%
        \par
    \fi
    \global\aenofloatfalse
}{}

\newmdenv[
    skipabove=1\topskip, skipbelow=0.6\topskip,
    innertopmargin=6pt, innerbottommargin=6pt,
    innerleftmargin=10pt, innerrightmargin=10pt,
    linewidth=0.9pt, roundcorner=4pt,
    leftline=true, rightline=true, topline=true, bottomline=true,
    linecolor=vizRule, backgroundcolor=vizBg,
    frametitlerule=true,
    frametitlefont=\sffamily\bfseries\color{vizRule},
    frametitlebackgroundcolor=vizBg,
    nobreak=true,
]{vizbox}

\newenvironment{vizcallout}[1][Visualization to build]{%
    \begin{vizbox}[frametitle={#1}]%
}{%
    \end{vizbox}%
}

% "What this rests on": the short list of conditions a statistical guarantee
% depends on, stated at the point where the guarantee is handed to the reader,
% each condition paired with the symptom you would see if it failed.
\definecolor{roBg}{HTML}{E9EEF5}      % pale slate background
\definecolor{roRule}{HTML}{1A476F}    % navy rule
\newmdenv[
    skipabove=1\topskip, skipbelow=0.6\topskip,
    innertopmargin=6pt, innerbottommargin=6pt,
    innerleftmargin=10pt, innerrightmargin=10pt,
    linewidth=0.9pt, roundcorner=4pt,
    linecolor=roRule, backgroundcolor=roBg,
    frametitlerule=true,
    frametitlefont=\sffamily\bfseries\color{roRule},
    frametitlebackgroundcolor=roBg,
    nobreak=true,
]{restsonbox}
\newenvironment{restson}[1][What this depends on]{%
    \begin{restsonbox}[frametitle={#1}]%
}{%
    \end{restsonbox}%
}

% Keep callout boxes on one page. A box taller than about a page opts out
% locally with [nobreak=false] (the roisim box in Chapter 8).
\renewmdenv[
    style=kaoboxstyle,
    backgroundcolor=RoyalBlue!25!White,
    frametitlebackgroundcolor=RoyalBlue!25!White,
    nobreak=true,
]{kaobox}

% Narrow text column (kaobook reserves a wide margin) means long monospace
% command names and compound terms occasionally overrun. emergencystretch lets
% TeX relax spacing as a last resort instead of producing overfull lines.
\setlength{\emergencystretch}{3em}

% Disable microtype font EXPANSION (keep protrusion). With the narrow column
% and unbreakable inline \texttt{} tokens, expansion hits its compression
% limit and whole paragraphs render visibly condensed, reading as a typeface
% change next to code blocks; it also garbles text extraction from the PDF.
\microtypesetup{expansion=false}

% --- Spacing and copy-paste fixes (Nov 2026) ---------------------------------
% 1. Real space characters in the PDF, so text copied from a PDF viewer keeps
%    its word breaks (pdfTeX otherwise positions words with no space glyph).
\input{glyphtounicode}\pdfgentounicode=1
\pdfinterwordspaceon
% 2. Straight quotes and backticks in verbatim, so Stata macro code such as
%    `today' copies from the book and runs.
\usepackage{upquote}
% 3. Tighter vertical space around verbatim code blocks: no extra list or
%    paragraph skip above and below beyond a small fixed gap.
\makeatletter
\preto{\@verbatim}{\topsep=6pt plus 1pt minus 1pt \partopsep=0pt \parskip=0pt}
\makeatother
% 4. Pages end at their natural height. With bottom justification, TeX
%    stretched the paragraph skip (0.5\baselineskip plus 0.5\baselineskip)
%    to fill each page, which opened large gaps after code blocks.
\raggedbottom
\AtBeginDocument{\setlength{\parskip}{0.5\baselineskip plus 1pt minus 0.5pt}}
% 5. Tighter section and subsection headings.
\RedeclareSectionCommand[beforeskip=-2.6ex plus -0.6ex minus -0.2ex,afterskip=1.0ex plus 0.2ex]{section}
\RedeclareSectionCommand[beforeskip=-2.0ex plus -0.5ex minus -0.2ex,afterskip=0.6ex plus 0.1ex]{subsection}

% Paths for figures
\graphicspath{{images/}}

% Display equations (align/aligned/cases); amssymb already loaded by kao.sty
\usepackage{amsmath}

% Flow diagrams (workflow pipelines, decision trees)
\usepackage{tikz}
\usetikzlibrary{arrows.meta, positioning, shapes.geometric, fit, backgrounds, calc}
% Reusable node styles for the book's pipeline diagrams
% kaobook turns every margin note into a float, so a chapter with figures and
% a run of margin notes can exhaust LaTeX's 18-slot float queue ("Too many
% unprocessed floats").  Ask the kernel for more slots rather than moving the
% figures around; this changes no layout.
\extrafloats{512}
% Full-page-width figures without the double-float machinery.  \begin{figure*}
% queues on LaTeX's separate double-float list, which a tufte-style book with
% hundreds of margin notes exhausts quickly ("Too many unprocessed floats").
% An ordinary float whose graphic spans the text block plus the margin column
% gives the same printed width and costs one ordinary float slot.
\newlength{\fullpagewidth}
\setlength{\fullpagewidth}{\dimexpr\textwidth+\marginparsep+\marginparwidth\relax}

\tikzset{
    stage/.style={rectangle, rounded corners=3pt, draw=vizRule, line width=0.9pt,
        fill=vizBg, text width=1.7cm, align=center, minimum height=1.0cm, font=\sffamily\footnotesize},
    substep/.style={rectangle, rounded corners=2pt, draw=aeRule, line width=0.6pt,
        fill=aeBg, align=left, font=\sffamily\scriptsize, inner sep=3pt},
    flow/.style={-{Stealth[length=2mm]}, line width=0.9pt, draw=aeRule},
    decision/.style={diamond, aspect=1.6, draw=vizRule, line width=0.8pt,
        fill=vizBg, align=center, inner sep=2pt, font=\sffamily\scriptsize},
    io/.style={trapezium, trapezium left angle=70, trapezium right angle=110,
        draw=aeRule, line width=0.7pt, fill=aeBg, align=center, font=\sffamily\scriptsize, inner sep=3pt},
    annot/.style={draw=none, fill=none, align=left, font=\itshape\scriptsize, text=black!70},
}

% Index & Glossary
\makeindex[columns=3, title=Alphabetical Index, intoc]
\makeglossaries

% Raise float ceiling: added text shifted anchors and jammed the 18-float queue (2026-07-23)
\extrafloats{100}
